\section{Referencias cruzadas}

\begin{frame}
	\frametitle{}

	Una de las razones para enumerar las diferentes partes en que se divide un documento \LaTeX,
	como por ejemplo, en capítulos y secciones, es poder dirigir al lector hacia una cierta
	parte del documento mediante las \alert{referencias cruzadas}.

	\LaTeX{} tiene un mecanismo automático para incluir las referencias cruzadas con ayuda de un fichero
	auxiliar \mintinline{latex}{.aux}, el cual es un subproducto de la última composición del documento.
	Por este motivo, se requiere procesar el documento \LaTeX{} al menos dos veces para obtener las referencias
	correctas.

	\begin{itemize}
		\item
		      \mintinline{latex}|\ref{<etiqueta>}|: Reemplaza en el documento compuesto el número de capítulo, apartado, figura, tabla o teorema
		      siempre que se haya incluido previamente \mintinline{latex}|\label{<etiqueta>}|.

		\item

		      \mintinline{latex}|\pageref{<etiqueta>}|: Reemplaza en el documento compuesto el número de página del documento.
	\end{itemize}
\end{frame}

\begin{frame}
	\frametitle{}
	\begin{columns}
		\begin{column}{.3\paperwidth}
			Las referencias bibliográficas se incluyen en un trabajo científico a fin de facilitar al lector
			la identificación precisa de una fuente de información.

			Los datos principales que se deben incluir en la referencia de un libro son
			\begin{enumerate}
				\item

				      Nombre del autor o los autores, o nombre del editor o de la institución que lo publica oficialmente, si no consta de ningún autor.

				\item

				      Identificación: título completo con el subtítulo, si lo hay.

				\item

				      Colaboradores de la edición: prologuistas, revisores, traductores, ilustradores.

				\item

				      Ordinal de la edición concreta utilizada, si no es la primera.

				\item

				      Datos de la publicación: lugar, editorial y año de publicación.

				\item

				      Colección y número.

				\item

				      Extensión (número de páginas, si se trata de un solo volumen, y número de volúmenes)
			\end{enumerate}
		\end{column}
		\begin{column}{.6\paperwidth}
			\begin{figure}[ht!]
				\centering
				\includegraphics[width=.2\paperwidth]{libro}
				\caption{Recuperado de~\url{https://fondoeduni.uni.edu.pe}.}
			\end{figure}
		\end{column}
	\end{columns}
\end{frame}

\begin{frame}
	\frametitle{¿Qué más por aprender?}
	\begin{itemize}
		\item \mintinline{latex}|fancyhdr|: Permite personalizar
		      fácilmente los encabezados y pies de página.

		\item \mintinline{latex}|makeidx|: Proporciona comandos para
		      producir índices.

		\item \mintinline{latex}|booktabs|: Proporciona comandos para
		      producir tablas con un formato atractivo para su documento:
	\end{itemize}
\end{frame}

\begin{frame}
	\begin{itemize}
		\item \mintinline{latex}|.tex|: archivo de entrada \LaTeX{}, se compila con \mintinline{latex}|lualatex|.

		\item \mintinline{latex}|.sty|: paquete de macros \LaTeX.
		      Cargue este paquete en su documento \LaTeX{} con el comando \mintinline{latex}|\usepackage|.

		\item \mintinline{latex}|.dtx|: archivo \TeX{} documentado.
		      Este es el principal formato de distribución para archivos
		      de estilo \LaTeX.
		      Si procesa un archivo \mintinline{latex}|.dtx|, obtendrá el
		      código de macro documentado del paquete \LaTeX{} contenido
		      en el archivo \mintinline{latex}|.dtx|.

		\item \mintinline{latex}|.ins|: El instalador de los archivos
		      incluidos en el archivo \mintinline{latex}|.dtx| correspondiente.
		      Si descarga un paquete \LaTeX{} de internet, normalmente
		      obtendrá un archivo \mintinline{latex}|.dtx| y uno
		      \mintinline{latex}|.ins|.
		      Ejecute \LaTeX en el archivo \mintinline{latex}|.ins| para
		      descomprimir el archivo \mintinline{latex}|.dtx|.

		\item \mintinline{latex}|.cls|: Los archivos de clase definen la
		      apariencia del documento.
		      Se seleccionan con el comando
		      \mintinline{latex}|\documentclass|.

		\item \mintinline{latex}|.fd|: Archivo de descripción de fuentes,
		      que informa a \LaTeX sobre nuevas fuentes.
	\end{itemize}
\end{frame}

\begin{frame}

\end{frame}
